%!TEX root = ../../main.tex

\chapter{Konzeption der experimentellen Evaluierung}
\label{ch:konzeption}

\section{Forschungsdesign und Hypothesen}
\label{sec:forschungsdesign}
% Übergeordnetes Design des Experiments, Hypothesen.

\section{Auswahl der Large Language Models}
\label{sec:modellauswahl}
% Auswahlkriterien und untersuchte Modelle (z.B. GPT-4, Claude, Llama).

\section{Definition der ETL-Aufgaben}
\label{sec:etl-aufgaben}

\subsection{Datenbereinigung}
\label{subsec:aufgabe-bereinigung}
% Behandlung fehlender Werte, Korrektur von Tippfehlern.

\subsection{Datentransformation}
\label{subsec:aufgabe-transformation}
% Formatvereinheitlichung, Typkonvertierung, Anreicherung.

\subsection{Duplikaterkennung}
\label{subsec:aufgabe-duplikate}
% Exakte und unscharfe Duplikate identifizieren.

\section{Prompting-Strategien}
\label{sec:prompting-strategien-konzept}
% Welche Strategien werden verglichen und warum.

\section{Evaluationskriterien}
\label{sec:evaluationskriterien}

\subsection{Genauigkeit}
\label{subsec:genauigkeit}

\subsection{Vollständigkeit}
\label{subsec:vollstaendigkeit}

\subsection{Konsistenz}
\label{subsec:konsistenz}

\subsection{Effizienz}
\label{subsec:effizienz}
% Laufzeit, Kosten, Token-Verbrauch.

\section{Vergleich mit klassischen Methoden als Baseline}
\label{sec:baseline}
% Definition der klassischen Vergleichsbasis.
